\documentclass[a4paper,12pt,twoside]{memoir}

% Castellano
\usepackage[spanish,es-tabla]{babel}
\selectlanguage{spanish}
\usepackage[utf8]{inputenc}
\usepackage[T1]{fontenc}
\usepackage{lmodern} % Scalable font
\usepackage{microtype}
\usepackage{placeins}

\RequirePackage{booktabs}
\RequirePackage[table]{xcolor}
\RequirePackage{xtab}
\RequirePackage{multirow}

% Links
\PassOptionsToPackage{hyphens}{url}\usepackage[colorlinks]{hyperref}
\hypersetup{
	allcolors = {red}
}

% Ecuaciones
\usepackage{amsmath}

% Rutas de fichero / paquete
\newcommand{\ruta}[1]{{\sffamily #1}}

% Párrafos
\nonzeroparskip

% Huérfanas y viudas
\widowpenalty100000
\clubpenalty100000

% Imágenes

% Comando para insertar una imagen en un lugar concreto.
% Los parámetros son:
% 1 --> Ruta absoluta/relativa de la figura
% 2 --> Texto a pie de figura
% 3 --> Tamaño en tanto por uno relativo al ancho de página
\usepackage{graphicx}
\newcommand{\imagen}[3]{
	\begin{figure}[!h]
		\centering
		\includegraphics[width=#3\textwidth]{#1}
		\caption{#2}\label{fig:#1}
	\end{figure}
	\FloatBarrier
}

% Comando para insertar una imagen sin posición.
% Los parámetros son:
% 1 --> Ruta absoluta/relativa de la figura
% 2 --> Texto a pie de figura
% 3 --> Tamaño en tanto por uno relativo al ancho de página
\newcommand{\imagenflotante}[3]{
	\begin{figure}
		\centering
		\includegraphics[width=#3\textwidth]{#1}
		\caption{#2}\label{fig:#1}
	\end{figure}
}

% El comando \figura nos permite insertar figuras comodamente, y utilizando
% siempre el mismo formato. Los parametros son:
% 1 --> Porcentaje del ancho de página que ocupará la figura (de 0 a 1)
% 2 --> Fichero de la imagen
% 3 --> Texto a pie de imagen
% 4 --> Etiqueta (label) para referencias
% 5 --> Opciones que queramos pasarle al \includegraphics
% 6 --> Opciones de posicionamiento a pasarle a \begin{figure}
\newcommand{\figuraConPosicion}[6]{%
  \setlength{\anchoFloat}{#1\textwidth}%
  \addtolength{\anchoFloat}{-4\fboxsep}%
  \setlength{\anchoFigura}{\anchoFloat}%
  \begin{figure}[#6]
    \begin{center}%
      \Ovalbox{%
        \begin{minipage}{\anchoFloat}%
          \begin{center}%
            \includegraphics[width=\anchoFigura,#5]{#2}%
            \caption{#3}%
            \label{#4}%
          \end{center}%
        \end{minipage}
      }%
    \end{center}%
  \end{figure}%
}

%
% Comando para incluir imágenes en formato apaisado (sin marco).
\newcommand{\figuraApaisadaSinMarco}[5]{%
  \begin{figure}%
    \begin{center}%
    \includegraphics[angle=90,height=#1\textheight,#5]{#2}%
    \caption{#3}%
    \label{#4}%
    \end{center}%
  \end{figure}%
}
% Para las tablas
\newcommand{\otoprule}{\midrule [\heavyrulewidth]}
%
% Nuevo comando para tablas pequeñas (menos de una página).
\newcommand{\tablaSmall}[5]{%
 \begin{table}
  \begin{center}
   \rowcolors {2}{gray!35}{}
   \begin{tabular}{#2}
    \toprule
    #4
    \otoprule
    #5
    \bottomrule
   \end{tabular}
   \caption{#1}
   \label{tabla:#3}
  \end{center}
 \end{table}
}

%
% Nuevo comando para tablas pequeñas (menos de una página).
\newcommand{\tablaSmallSinColores}[5]{%
 \begin{table}[H]
  \begin{center}
   \begin{tabular}{#2}
    \toprule
    #4
    \otoprule
    #5
    \bottomrule
   \end{tabular}
   \caption{#1}
   \label{tabla:#3}
  \end{center}
 \end{table}
}

\newcommand{\tablaApaisadaSmall}[5]{%
\begin{landscape}
  \begin{table}
   \begin{center}
    \rowcolors {2}{gray!35}{}
    \begin{tabular}{#2}
     \toprule
     #4
     \otoprule
     #5
     \bottomrule
    \end{tabular}
    \caption{#1}
    \label{tabla:#3}
   \end{center}
  \end{table}
\end{landscape}
}

%
% Nuevo comando para tablas grandes con cabecera y filas alternas coloreadas en gris.
\newcommand{\tabla}[6]{%
  \begin{center}
    \tablefirsthead{
      \toprule
      #5
      \otoprule
    }
    \tablehead{
      \multicolumn{#3}{l}{\small\sl continúa desde la página anterior}\\
      \toprule
      #5
      \otoprule
    }
    \tabletail{
      \hline
      \multicolumn{#3}{r}{\small\sl continúa en la página siguiente}\\
    }
    \tablelasttail{
      \hline
    }
    \bottomcaption{#1}
    \rowcolors {2}{gray!35}{}
    \begin{xtabular}{#2}
      #6
      \bottomrule
    \end{xtabular}
    \label{tabla:#4}
  \end{center}
}

%
% Nuevo comando para tablas grandes con cabecera.
\newcommand{\tablaSinColores}[6]{%
  \begin{center}
    \tablefirsthead{
      \toprule
      #5
      \otoprule
    }
    \tablehead{
      \multicolumn{#3}{l}{\small\sl continúa desde la página anterior}\\
      \toprule
      #5
      \otoprule
    }
    \tabletail{
      \hline
      \multicolumn{#3}{r}{\small\sl continúa en la página siguiente}\\
    }
    \tablelasttail{
      \hline
    }
    \bottomcaption{#1}
    \begin{xtabular}{#2}
      #6
      \bottomrule
    \end{xtabular}
    \label{tabla:#4}
  \end{center}
}

%
% Nuevo comando para tablas grandes sin cabecera.
\newcommand{\tablaSinCabecera}[5]{%
  \begin{center}
    \tablefirsthead{
      \toprule
    }
    \tablehead{
      \multicolumn{#3}{l}{\small\sl continúa desde la página anterior}\\
      \hline
    }
    \tabletail{
      \hline
      \multicolumn{#3}{r}{\small\sl continúa en la página siguiente}\\
    }
    \tablelasttail{
      \hline
    }
    \bottomcaption{#1}
  \begin{xtabular}{#2}
    #5
   \bottomrule
  \end{xtabular}
  \label{tabla:#4}
  \end{center}
}



\definecolor{cgoLight}{HTML}{EEEEEE}
\definecolor{cgoExtralight}{HTML}{FFFFFF}

%
% Nuevo comando para tablas grandes sin cabecera.
\newcommand{\tablaSinCabeceraConBandas}[5]{%
  \begin{center}
    \tablefirsthead{
      \toprule
    }
    \tablehead{
      \multicolumn{#3}{l}{\small\sl continúa desde la página anterior}\\
      \hline
    }
    \tabletail{
      \hline
      \multicolumn{#3}{r}{\small\sl continúa en la página siguiente}\\
    }
    \tablelasttail{
      \hline
    }
    \bottomcaption{#1}
    \rowcolors[]{1}{cgoExtralight}{cgoLight}

  \begin{xtabular}{#2}
    #5
   \bottomrule
  \end{xtabular}
  \label{tabla:#4}
  \end{center}
}



\graphicspath{ {./img/} }

% Capítulos
\chapterstyle{bianchi}
\newcommand{\capitulo}[2]{
	\setcounter{chapter}{#1}
	\setcounter{section}{0}
	\setcounter{figure}{0}
	\setcounter{table}{0}
	\chapter*{\thechapter.\enskip #2}
	\addcontentsline{toc}{chapter}{\thechapter.\enskip #2}
	\markboth{#2}{#2}
}

% Apéndices
\renewcommand{\appendixname}{Apéndice}
\renewcommand*\cftappendixname{\appendixname}

\newcommand{\apendice}[1]{
	%\renewcommand{\thechapter}{A}
	\chapter{#1}
}

\renewcommand*\cftappendixname{\appendixname\ }

% Formato de portada
\makeatletter
\usepackage{xcolor}
\newcommand{\tutor}[1]{\def\@tutor{#1}}
\newcommand{\course}[1]{\def\@course{#1}}
\definecolor{cpardoBox}{HTML}{E6E6FF}
\def\maketitle{
  \null
  \thispagestyle{empty}
  % Cabecera ----------------
\noindent\includegraphics[width=\textwidth]{cabecera}\vspace{1cm}%
  \vfill
  % Título proyecto y escudo informática ----------------
  \colorbox{cpardoBox}{%
    \begin{minipage}{.8\textwidth}
      \vspace{.5cm}\Large
      \begin{center}
      \textbf{TFG del Grado en Ingeniería Informática}\vspace{.6cm}\\
      \textbf{\LARGE\@title{}}
      \end{center}
      \vspace{.2cm}
    \end{minipage}

  }%
  \hfill\begin{minipage}{.20\textwidth}
    \includegraphics[width=\textwidth]{escudoInfor}
  \end{minipage}
  \vfill
  % Datos de alumno, curso y tutores ------------------
  \begin{center}%
  {%
    \noindent\LARGE
    Presentado por \@author{}\\ 
    en Universidad de Burgos --- \@date{}\\
    Tutor: \@tutor{}\\
  }%
  \end{center}%
  \null
  \cleardoublepage
  }
\makeatother

\newcommand{\nombre}{Jorge Martínez Martín} %%% cambio de comando

% Datos de portada
\title{Clasificación de grado de Parkinson mediante aprendizaje semisupervisado}
\author{\nombre}
\tutor{Álvar Arnaiz González}
\date{\today}

\begin{document}

\maketitle


\newpage\null\thispagestyle{empty}\newpage


%%%%%%%%%%%%%%%%%%%%%%%%%%%%%%%%%%%%%%%%%%%%%%%%%%%%%%%%%%%%%%%%%%%%%%%%%%%%%%%%%%%%%%%%
\thispagestyle{empty}


\noindent\includegraphics[width=\textwidth]{cabecera}\vspace{1cm}

\noindent D. Álvar Arnaiz González, profesor del departamento de Ingeniería Informática, área de Lenguajes y Sistemas Informáticos.

\noindent Expone:

\noindent Que el alumno D. \nombre, con DNI 71482657Z, ha realizado el Trabajo final de Grado en Ingeniería Informática titulado título de TFG. 

\noindent Y que dicho trabajo ha sido realizado por el alumno bajo la dirección del que suscribe, en virtud de lo cual se autoriza su presentación y defensa.

\begin{center} %\large
En Burgos, {\large \today}
\end{center}

\vfill\vfill\vfill

% Author and supervisor
\begin{minipage}{0.45\textwidth}
\begin{flushleft} %\large
Vº. Bº. del Tutor:\\[2cm]
D. nombre tutor
\end{flushleft}
\end{minipage}
\hfill
\begin{minipage}{0.45\textwidth}
\begin{flushleft} %\large
Vº. Bº. del co-tutor:\\[2cm]
D. nombre co-tutor
\end{flushleft}
\end{minipage}
\hfill

\vfill

% para casos con solo un tutor comentar lo anterior
% y descomentar lo siguiente
%Vº. Bº. del Tutor:\\[2cm]
%D. nombre tutor


\newpage\null\thispagestyle{empty}\newpage




\frontmatter

% Abstract en castellano
\renewcommand*\abstractname{Resumen}
\begin{abstract}
La enfermedad del Parkinson afecta cada vez a más personas. Pese a ser una enfermedad incurable poder detectarla y ver su progreso es crucial para adecuar los tratamientos a cada persona.

\end{abstract}

\renewcommand*\abstractname{Descriptores}
\begin{abstract}
Palabras separadas por comas que identifiquen el contenido del proyecto Ej: servidor web, buscador de vuelos, android \ldots
\end{abstract}

\clearpage

% Abstract en inglés
\renewcommand*\abstractname{Abstract}
\begin{abstract}
A \textbf{brief} presentation of the topic addressed in the project.
\end{abstract}

\renewcommand*\abstractname{Keywords}
\begin{abstract}
keywords separated by commas.
\end{abstract}

\clearpage

% Indices
\tableofcontents

\clearpage

\listoffigures

\clearpage

\listoftables
\clearpage

\mainmatter
\capitulo{1}{Introducción}

Descripción del contenido del trabajo y del estructura de la memoria y del resto de materiales entregados.

\capitulo{2}{Objetivos del proyecto}

El principal objetivo de este proyecto es la creación de una aplicación web que permita, mediante la evaluación del grado de bradicenesia utilizando vídeos de la prueba del <<rapid finger tapping test>> y el uso de inteligencia artificial realizar dicha evaluación. Además, esta aplicación debe poder ser usada por el usuario promedio.


\section{Objetivos técnicos}

Por ello, para poder alcanzar estos objetivos generales, se plantean los siguientes objetivos técnicos:
\begin{itemize}
\item \textbf{Revisar de trabajos relacionados}: para comprender mejor el campo que se esta investigando y obtener conocimientos relacionados con el tema.
\item \textbf{Desarrollar el modelo de aprendizaje supervisado}: entrenar, evaluar y seleccionar el modelo que mejor resultados ofrezca.
\item \textbf{Evaluar la calidad de la solución}: comprobando si los resultados que nos ofrece son coherentes.
\item \textbf{Crear de una aplicación \textit{web}}: que sea fácil de usar e intuitiva para el usuario.
\item \textbf{Familiarizarse con nuevas tecnologías}: como programación \textit{web} o Docker.
\item \textbf{Documentar el proyecto}: de manera clara y concisa, condensando toda la información relevante en un formato accesible y fácil de entender.
\end{itemize}

\section{Objetivos de \textit{software}}

Se han propuesto los siguiente objetivos de desarrollo de \textit{software}.
\begin{itemize}
\item \textbf{Crear una interfaz de usuario accesible}: para que el usuario medio pueda usarla sin dificultad.
\item \textbf{Crear un sistema de gestión de datos persistente}: como bases de datos SQL.
\item \textbf{Asegurar la calidad del código}: mediante herramientas de análisis de código como SonarCloud.
\item \textbf{Deplegar en plataformas}: mediante el uso de herramientas como Docker, permitiendo además la familiarización con estas tecnologías.
\item \textbf{Implementar trabajos anteriores}: para realizar la extracción de características y obtener pautas para el correcto desarrollo.
\end{itemize}
\capitulo{3}{Conceptos teóricos}

En este capítulo se definirán algunos conceptos teóricos para facilitar la comprensión de este proyecto.

\section{Enfermedad del Parkinson}

La EP es una  enfermedad neurodegenerativa multisistémica progresiva que cada afecta principalmente a a gente de avanzada edad~\cite{pdsymptoms}.  
Entre los síntomas principales de esta enfermedad podemos encontrar:
Perdida significativa de parte de las células productoras de dopamina lo que produce que aparezcan mucho antes los síntomas relacionados con el movimiento que presenta la EP. Otro de los síntomas que podemos encontrar en la EP es la bradicinesia es cual se caracteriza por la lentitud al realizar movimientos voluntarios así como la ralentización y decremento de amplitud a la hora de realizar movimientos repetitivos.
El diagnóstico de la EP se basa en criterios específicos. Los síntomas iniciales incluyen lentitud de movimientos, junto con rigidez muscular, temblores o problemas de equilibrio. Los controles posteriores descartan otras posibles causas y confirman la presencia de factores que sugieren claramente la EP.
La EP suele empezar en un lado del cuerpo y extenderse a lo largo de unos años. Los síntomas incluyen postura encorvada, rigidez, letra más pequeña y marcha arrastrando los pies. Los temblores son frecuentes.
Las alteraciones de la marcha (como pasos indecisos o congelación repentina) se vuelven habituales. La pérdida de equilibrio es un problema importante, que aumenta el riesgo de caídas y lesiones.

\subsection{Rapid finger tapping test}

El test de golpeteo rápido de los dedos, también conocido como <<rapid finger tapping test>> o RFTT es un procedimiento por el cual el paciente realiza golpeteos de manera repetitiva durante un periodo de entre 10 y 15 segundos en los cuales ha de intentar generar la mayor amplitud posible entre el dedo indice y el pulgar sin bajar la frecuencia a la que lo realiza. Este test es usado comúnmente para el diagnostico de personas con la EP ya que como se explicó anteriormente uno de los síntomas que presenta esta enfermedad el la bradicisnesia el cual produce que una persona con EP al realizar este test muestre deteriodo ya sea en la amplitud o en la frecuencia según avanza la prueba. Este test será el que se utilizará en este proyecto para determinar el grado de EP en el que se encuentra el paciente.

\subsection{Unified parkinson disease rating scale}

La escala de evaluación unificada de la EP o UPDRS por sus siglas en inglés es una herramienta creada por la Movement Disorder Society que permite, mediante la evaluación de diversos parámetros del la EP, medir la gravedad de la enfermedad. En ella se miden diversos aspectos de las experiencias tanto motoras como no motoras de la vida diaria. Para el proyecto utilizaremos solo algunas de las medidas que se utilizan para la clasificación del estado de la persona ya que son los datos que se extraerán de los vídeos proporcionados por los pacientes.


\section{Aprendizaje automático}

El aprendiza automático es una rama de la inteligencia artificial centrada en el desarrollo de métodos y algoritmos para que el computador de manera autónoma sea capaz de, mediante la experiencia y el procesamiento de datos, mejorar en esa tarea. De esta manera los modelos realizaran predicciones cada vez más precisas.

Dentro de un <<dataset>> podemos encontrar lo que se conocen como instancias. Una instancia es cada fila del <<dataset>> caracterizadas por atributos que pueden ser tanto categóricos (nombres, colores, categorías,etc) como valores numéricos (números).

Dentro del aprendizaje automático podemos distinguir tres tipos principales:
\begin{itemize}
\item Aprendizaje supervisado: a este aprendizaje se le proporciona un con junto de datos completo con ejemplos etiquetados, y es mediante estos ejemplos con los que aprende a clasificar nuevos datos o realizar predicciones. Se puede dividir en regresión (predicción de datos numéricos) o clasificación (predicción de datos categóricos)
\item Aprendizaje no supervisado: este aprendizaje a diferencia del supervisado se le entregan los datos sin clasificar de manera que tiene que ser él el que encuentre las relaciones. Se puede dividir en: clustering (agrupación de datos por su similitud) y reducción dimensional (disminuye el número de variables en un conjunto de datos sin perder información importante).
\item Aprendizaje semi-supervisado: este aprendizaje es una combinación de los dos aprendizajes anteriores ya que se le proporcionan dos conjuntos de  datos. En el primer conjuntos los datos están etiquetados, este conjunto es el conocido como entrenamiento (training) que es con el cual el modelo se entrena para encontrar la relación entre los datos. En el segundo conjunto los datos no están etiquetados, este conjunto es conocido como test que es el que permite al modelo comparar como de precisas son las predicciones que se realizan ya que, pese a que al modelo se le introduzcan los datos sin etiquetar, el <<dataset>> si contiene las etiquetas de estas instancias. Al igual que el aprendizaje supervisado su objetivo principal es la creación de clasificadores o regresores.
\end{itemize}

Un clasificador se compone principalmente de dos fases: la fase de entrenamiento y la fase de predicción. Durante la primera fase se le introducen datos al clasificador con una etiqueta y el objetivo de este es encontrar una hipótesis que se cumpla para cualquier instancia incluso si no se encuentra dentro de los datos de entrenamiento. Para ello suponemos que la  hipótesis que mejor se adapte a las instancias que no se encuentran en la fase de entrenamiento, es la que mejor se adapte al conjunto de entrenamiento. Es decir, el clasificador <<aprende>> del conjunto de entrenamiento. 








\capitulo{4}{Técnicas y herramientas}

En este capítulo se muestran las diferentes técnicas y herramientas que se han utilizado para el desarrollo del proyecto.


\section{Técnicas}

En esta sección se muestran las técnicas principales empleadas para el desarrollo del proyecto.

\subsection{Scrum}

Para el desarrollo de este proyecto ha sido utilizada la metodología ágil conocida como Scrum.

La metodología Scrum trata de fraccionar la duración de un proyecto en lo que se conoce como <<sprints>>, cuya duración varía de una a dos semanas.
El objetivo de los <<sprints>> es decidir que parte del proyecto se va a desarrollar durante ese periodo de tiempo realizando revisiones diarias para ver como avanza asi como una revisión de <<sprint>> en el que hace una valoración general del <<sprint>> y se decide en que va a consistir el siguiente.

\section{Herramientas}

En esta sección se describen las herramientas que se han utilizado durante la realización del proyecto.

\subsection{Portales}

Usados principalmente para el seguimiento de la metodología Scrum, creación de diagramas y bocetos y control de versiones:
\begin{itemize}
\item \textbf{Zube}: portal dedicado a la gestión de proyectos <<software>>~\cite{zubeHome}. Permite visualizar proyectos, sprints,
gráficos propios de Scrum y crear tableros.
\item \textbf{GitHub}: es una plataforma web que proporciona una interfaz gráfica para el control de versiones utilizando Git~\cite{gitbubHome}. Además, GitHub permite a los usuarios revisar versiones anteriores de proyectos y realizar un seguimiento de los cambios a lo largo del tiempo.
\item \textbf{DrawIO}: es una herramienta que permite la creación de diagramas entidad-relación, diagramas de flujo, casos de uso, etc~\cite{drawioHome}. Esta herramienta dispone además de los símbolos UML necesarios para todas las funcionalidades.
\end{itemize}

\subsection{Librerías}
Principalmente de Python:
\begin{itemize}
\item \textbf{Scikit-Learn}: esta librería ofrece una amplia variedad de algoritmos de aprendizaje supervisados, semisupervisados y no supervisados~\cite{sklearnHome}. Además posee una gran compatibilidad con otras librerías como así como una extensa documentación que facilitan su implementación.3.
\item \textbf{TSFresh}: es una biblioteca de Python, para la extracción y representación de características de series temporales~\cite{tsfreshHome}. Esta herramienta agiliza el proceso de extracción de características a partir de datos de series temporales para apoyar las tareas de aprendizaje automático. Ofrece una selección de características definidas que abarcan diferentes medidas estadísticas, como análisis de tendencias, patrones estacionales, métricas de correlación y evaluaciones de complejidad.
\end{itemize}

\subsection{Entorno de desarrollo}
Incluye el entorno integrado de desarrollo (IDE), lenguaje y otros programas usados.
\begin{itemize}
\item \textbf{Visual Studio Code}: editor y depurador de código, así como algunas de sus extensiones más importantes.
\item \textbf{Python}: lenguaje de programación de alto nivel ampliamente usado. Ha sido escogido para realizar este proyecto debido a que posee una gran variedad de paquetes y librerías de Machine Learning~\cite{python}.
\item \textbf{Git}: es un software de control de versiones distribuido, utilizado para la control de versiones en proyectos software.
\item \textbf{Github Copilot}: es un asistente de programación con inteligencia artificial que te permite realizar código de forma más rápida mediante la realización de consultas y las sugerencias que provee.
\end{itemize}
\subsection{Desarrollo web}
A continuación se exponen tanto librerías como recursos para el desarrollo del apartado web:

\begin{itemize}
\item \textbf{Boostrap}: es un <<framework>> utilizado para la creación de páginas web~\cite{boostrap}. En esencia es un librería de código que simplifica el proceso de desarrollo de una página web. Unas de sus principales características son:
\begin{itemize}
\item \textbf{Diseño <<responsive>>}: Boostrap permite de una forma muy sencilla implementar la posibilidad de que la página web con todos sus componentes se adapten automáticamente en tiempo real a cualquier pantalla incluso si se modifica el tamaño de la ventana.
\item \textbf{Documentación y comunidad}: dado que Boostrap es un framework muy conocido, tanto la documentación del framework como la comunidad que ayuda a desarrollarlo facilitan mucho la posibilidad de aprenderlo y poder encontrar soluciones especificas a nuestro diseño en poco tiempo.
\item \textbf{Enfoque móvil}: Boostrap esta diseñado teniendo en cuenta también el diseño en dispositivos móviles, lo que ayuda al desarrollo de la web para multiplataformas.
\end{itemize}
\item \textbf{Flask}: es un <<framework>> escrito en Python  que simplifica y facilita la creación de aplicaciones web mediante el uso del patrón Modelo-Vista-Controlador~\cite{flask}. Dado que es considerado un micro <<framework>> no viene con todas las funcionalidades por defecto; sin embargo, existen librerías compatibles con Flask que permiten añadir todas las funcionalidades que se necesiten. Entre ellas se han usado:
\begin{itemize}
\item \textbf{Werkzeug}: es un conjunto de librerías de Python usados para el desarrollo web~\cite{werkzeug}. Se ha utilizado Flask junto a Werkzeug ya que este contiene una librería de cifrado que se ha utilizado durante el desarrollo del proyecto.
\item \textbf{SQLAlchemy}: dado que Flask no puede hacer uso de forma directa de las bases de datos hace uso de la librería SQLAlchemy~\cite{sqlalchemy}. Esta librería permite la creación de bases de datos mediante la creación de clases en Python, definir el esquema y la relación entre las tablas y hacer consultas mediante métodos y objetos en vez de consultas SQL sin formato.
\item \textbf{Flask-Login}: es una extensión de Flask que simplifica añadir la funcionalidad de iniciar sesión, administrar sesión y cerrar sesión~\cite{flasklogin}. Permite además otorgar y restringir accesos de distintas secciones de la aplicación así como cargar datos relacionados con la sesión actual.
\end{itemize}
\item \textbf{Jinja}: es un motor de plantillas rápido, seguro y fácil de usar para Python~\cite{jinja}. Permite la modificación de partes de la plantilla de la web mediante el uso de código como la agregación de contenido dependiendo del valor de variables mediante bucles o condicionales.
\item \textbf{Font Awesome}: es una herramienta usada para el diseño y desarrollo de páginas web~\cite{fontawesome}. En vez de usar archivos de imágenes para iconos utiliza fuentes. Esto ofrece diversas ventajas:
\begin{itemize}
\item Escalabilidad: los iconos se escalan perfectamente a cualquier tamaño sin perder calidad.
\item Colores personalizables: permite modificar el color de los iconos usando solo CSS.
Para poder utilizar Font Awesome con Flask se ha utilizado la librería Flask-FontAwesome que permite su implementación~\cite{flaskfontawesome}.
\end{itemize}
\end{itemize}
\capitulo{5}{Aspectos relevantes del desarrollo del proyecto}

En capítulo de la memoria pretende mostrar y resumir los aspectos relevantes durante el desarrollo del proyecto. En el se incluyen decisiones tomadas durante el desarrollo, modelado del problema y resultados de los experimentos.

En la subsección~\ref{Fase de experimentación} se explica y desarrolla la fase de experimentación con modelos. Para ello se expone como se ha realizado la extracción de características, que algoritmos se han utilizado durante la experimentación, como se ha realizado el tratamiento de los datos y como se ha realizado el entrenamiento y la evaluación de los modelos.

En la subsección~\ref{Resultados} se exponen los resultados mediante gráficas y se expone el análisis realizado para seleccionar el modelo que mejor se adecúa al problema.

Por último en la subsección~\ref{Fase de desarrollo de la aplicación} se exponen detalles de la fase de desarrollo entre los que se incluyen el motivo de las tecnologías utilizadas y problemas durante el desarrollo.

\section{Fase de experimentación}
\label{Fase de experimentación}
\subsection{Extracción de características mediante un Trabajo de Fin de Grado previo}

Para la extracción de características se hizo uso de un apartado del Trabajo de Fin de Grado de Catalin llamado PADDEL~\cite{paddelRepo}. Este TFG consiste en detectar mediante la bradicenesia si una persona tiene o no Parkinson. Para ello extrae las características de vídeos haciendo el RFTT mediante visión artificial y, aplicando técnicas de minería de datos, es capaz de predecir si una persona tiene o no la enfermedad de Parkinson.

Realizando un estudio del TFG de Catalin se pudo extraer y utilizar correctamente el código que permitía la extracción de características. Este código, tras algunas modificaciones, ha sido implementado en el proyecto.

Las características que se extraen de los vídeos, se pueden dividir en tres grupos:
\begin{itemize}
\item Atributos iniciales: son los atributos que de los que se dispone mediante el titulo del vídeo como son la edad, el sexo o la mano dominante.
\item Atributos extraídos del vídeo: son los atributos que extrae mediante series temporales como son la velocidad, la amplitud o el número de golpes.
\item Atributos extraídos de vídeo con TSFresh: Son los 794 atributos
del vídeo extraídos mediante TSFresh.
\end{itemize}


\subsection{Análisis y tratamiento de los datos}
Una vez se extrajeron las características de los datos se pudo llevar a cabo un análisis de los datos.
En el conjunto de datos disponible en la fase de experimentación se disponía de 222 vídeos. 
Cada video disponía de los siguientes atributos:
\begin{itemize}
\item Sexo de la persona.
\item Edad de la persona.
\item Mano dominante de la persona.
\item Mano que aparece en el vídeo.
\item Fecha en la que se ha tomado.
\item Valoración del grado de lentitud del vídeo. Primera variable objetivo.
\item Valoración del grado de amplitud del vídeo. Segunda variable objetivo.
\end{itemize}
Las valoraciones fueron facilitadas por dos neurólogas del HUBU. Estas valoraciones diferían  entre sí de manera considerable y es por ello se decidió realizar el experimento utilizando solo las clasificaciones de una de las dos. 

Una vez analizados los datos y extraídas las características se hicieron varías comprobaciones. Primero se comprobó que todos los vídeos estuviesen correctamente valorados. Este no fue el caso ya que 2 de los videos de los que se disponía no estaban valorados es por ello que, debido a que el tipo de aprendizaje utilizado es supervisado, se decidió que no se incluirían en la fase de experimentación.

Tras eliminar los vídeos que no fueron valorados, se comprobó si las clases estaban desbalanceadas. Al realizar esta comprobación se vio que había un desbalanceado muy elevado, habiendo clases que solo tenían menos de 8 instancias. Dado que se pretendía utilizar algoritmos como KNN o técnicas de remuestreo como \textit{Random Oversampling} se tomo la decisión de eliminar aquellas clases que tuviesen menos de 8 instancias. Las clases que se eliminaron fueron lentitud 3 y 4 así como amplitud 4.
\imagen{./graphs/class_graph1.png}{Número de instancias por clase antes de realizar el tratamiento de datos.}{}
Además de lo anteriormente mencionado se comprobó si existía alguna característica de los vídeos que fuese nula. Al comprobarlo se observo que en algunas todas había una característica que era nula por lo que se decidió eliminar esa característica nula de todas las instancias ya que no aportaba ninguna información adicional. Una de esas características era la edad, que solo fue nula en vídeos de control, es decir vídeos de personas sin la enfermedad de Parkinson y dado al elevado número de vídeos de control, no supuso una gran perdida de información. Posteriormente se eliminaron aquellos vídeos cuya duración no era superior a 15 segundos ya que, según aparece en la documentación de PADDEL, no se realizaba correctamente la extracción de características si el vídeo no cumplía ese requisito.
\imagen{./graphs/class_graph1.png}{Número de instancias por clase tras realizar el tratamiento de datos.}{}
Tras eliminar las instancias anteriormente mencionadas, se realizó un tratamiento de datos en el que las característica categóricas se transformaron en numéricas para poder entrenar correctamente los modelos.

\subsection{Elección del modelo}

Una vez se trataron los datos correctamente, se realizó una selección de los modelos a evaluar. Para ello se eligieron aquellos modelos que mejor funcionaron en las pruebas de PADDEL así como los que obtuvieron mejores resultados en el artículo <<Computer Vision for Parkinson’s Disease Evaluation: A Survey on Finger Tapping>>~\cite{AmoSalas2024}. Los algoritmos seleccionados fueron KNN, \textit{Random Forest}, XGBoost y SVM.

Para la comparación de modelos, se utilizo \textit{GridSearchCV}, una técnica de validación cruzada que nos permite evaluar los mejores hiperparámetros para un modelo. Para ello el modelo se entrena utilizando una parte del \textit{dataset} como entrenamiento y la restante como prueba. Para estas pruebas se decidió utilizar 320 características de las 807 resultantes tras el tratamiento de los datos. Este dato fue extraído de la documentación del proyecto de PADDEL~\cite{paddelRepo} y probado durante los experimentos viendo que era el que mejor resultados reportaba (ver comparativa en las figuras~\ref{fig:./graphs/comparativeLentSmote.png} y~\ref{fig:./graphs/comparativeLentSmote320.png}). Es importante destacar que para entrenar cada modelos, las estadísticas que se seleccionan difieren ya que no tienen las mismas correlaciones.
\imagen{./graphs/comparativeLentSmote.png}{Resultados de modelos con todas las características y SMOTE}{}
Durante la fase de entrenamiento se hizo uso tanto de \textit{Random Over-Sampling} (ROS) como de \textit{Synthetic Minority Over-sampling Technique} (SMOTE). Estas técnicas son esenciales en problemas de desequilibrio de clases, ya que permiten entrenar el modelo haciendo que todas las clases tengan el mismo número de instancias y evitar que el modelo ignore las clases minoritarias y pueda predecirlas adecuadamente. Al utilizar este tipo de técnicas hay que tener en cuenta que solo se ha de hacer remuestreo sobre los datos de entrenamiento ya que los datos de test no se deben modificar. 

Esto supuso un problema ya que internamente \textit{GridSearchCV} no te permite separar la parte de entrenamiento y test y hacer el \textit{Over-Sampling} correctamente. Es por ello que se decidió hacer de manera externa al entrenamiento del modelo una validación cruzada estratificada dividida en 5 grupos. En ella 4 de los 5 grupos eran utilizados para entrenar el modelo y el restante se utilizaba para evaluar el rendimiento del modelo con los datos no modificados.

Dado que este problema se componía de dos variables objetivo, se opto por realizar modelos independientes para cada clase de manera que cada modelo predijese una variable. Es por ello que, durante el entrenamiento de cada una de las variable objetivo, la variable que no se predecía se omitía y no se introducía durante los entrenamientos.

\section{Resultados}
\label{Resultados}

Tras realizar el tratamiento de los datos y el entrenamiento de los modelos se obtuvieron 20 grupos de evaluaciones para cada modelo. Estos grupos están compuestos de las siguientes evaluaciones:
\begin{itemize}
\item \textit{Accuracy}
\item \textit{F-score}
\item \textit{G-mean}
\end{itemize}
Cada modelo fue evaluado con cinco veces, una por cada grupo dentro de las divisiones de la validación cruzada y dentro de cada grupo fue evaluado utilizando ROS y SMOTE.

Para realizar acabo estas evaluaciones, se utilizó tanto en la medida F-score como en la \textit{G-mean} el parámetro de \textit{average} en \textit{weighted} debido al gran desbalanceo entre clases ya que, si se hubiese indicado \textit{macro}, debido a que internamente realiza una media aritmética entre cada \textit{F-score} de cada clase, no se habría tenido en cuenta el desbalanceo.

Para obtener una información más general sobre los modelos y sus resultados se decidió que las gráficas que representaran los valores fuesen diagramas de cajas y bigotes, lo cual permitió un mejor visión general del cada modelo.

\subsection{Lentitud}
Para el caso de la lentitud, se hicieron de forma paralela la evaluación del modelo para la lentitud utilizando SMOTE y ROS (ver figuras~\ref{fig:./graphs/comparativeLentSmote320.png} y~\ref{fig:./graphs/comparativeLentRos320.png} respectivamente).
\imagen{./graphs/comparativeLentSmote320.png}{Resultados de modelos con 320 características y SMOTE.}{}
\imagen{./graphs/comparativeLentRos320.png}{Resultados de modelos con 320 características y ROS.}{}
Para hacer las comparativas entre los modelos, se omitió la \textit{Accuracy} ya que, en una clasificación multiclase, no es una buena medida para indicar como de bueno es el modelo. Es por esto que las medidas utilizadas para la comparación entre los modelos fueron el \textit{F-score} y \textit{G-mean}. No obstante, no se descartó la medida de \textit{Accuracy} durante la experimentación ya que proveía una vista general del rendimiento de los modelos frente a un problema desconocido.

En términos generales se puede observar que SMOTE tiene mejor desempeño que ROS. Tras descartar ROS, la comparación entre los modelos muestra que el algoritmo de KNN tiene un desempeño peor comparado con respecto al resto de modelos. Es por ello que queda descartado como modelo candidato. Entre los modelos que quedan \textit{XGBoost} y \textit{Random Forest} destacan sobre el SVM, descartandolo así como modelo candidato. Por último, entre los dos modelos restantes, \textit{XGBoost} pese a tener una mediana peor, posee un valor máximo que destaca en gran medida con respecto a \textit{Random Forest}. Es por ello que se decidió utilizar \textit{XGBoost} como modelo para la lentitud.

Dado que este diagrama es la representación de cinco iteraciones, los mejores hiperparámetros son aquellos que proporcionen el valor máximo. Pese a tener dos valores máximos para dos medidas distintas, los valores pertenecen a la misma instancia por lo que solo tenemos un grupo de hiperparámetros.

Para el caso de la lentitud los hiperparámetros fueron:
\begin{itemize}
\item \textit{grow\_policy}: \textit{depthwise}.
\item \textit{learning\_rate}: 0,1.
\item \textit{Number of estimators}: 200.
\end{itemize}
Este modelo con los hiperparámetros indicados obtuvo un \textit{F-score} de 0,69 y un \textit{G-mean} de 0,71 siendo estos los mejores valores reportados entre todos los modelos.

\subsection{Amplitud}

Al igual que en el caso de la lentitud,  se hicieron de forma paralela la evaluación del modelo para la lentitud utilizando ROS y SMOTE (ver figuras~\ref{fig:./graphs/comparativeAmpRos320.png} y~\ref{fig:./graphs/comparativeAmpSmote320.png} respectivamente).
\imagen{./graphs/comparativeAmpRos320.png}{Resultados de modelos con 320 características y ROS.}{}
\imagen{./graphs/comparativeAmpSmote320.png}{Resultados de modelos con 320 características y SMOTE.}{}

En el caso de la amplitud, la diferencia entre ROS y SMOTE no es tan pronunciada; sin embargo, podemos observar una diferencia notoria en el caso de \textit{XGBoost} el cual se ve claramente beneficiado por SMOTE. El resto de modelos, aunque no de forma tan notoria también se ven beneficiados por SMOTE. Es por ello que se decide tomar SMOTE como método de remuestreo como se hizo para la lentitud. 

Parecido a lo que ocurría con la lentitud KNN tiene peor desempeño en el caso de \textit{F-sore} y SVM, \textit{Random Forest} y \textit{XGBoost} tienen resultados muy parecidos. En el caso de \textit{G-mean} tanto SVM como \textit{Random Forest} presentan un peor rendimiento, mientras que los dos modelos restantes obtienen mejores resultados.

Dentro de las comparativas de SMOTE, vuelve a destacar \textit{XGBoost} que, en esta ocasión, posee un mejor rendimiento general mejor que el resto de modelos a diferencia de lo que ocurría con la lentitud que se eligió por el caso máximo. 

Como ya se explicó anteriormente al ser diagramas de cajas y bigotes, se elegirá el máximo ya que reporta los mejores rendimientos.

Para el caso de la amplitud los hiperparámetros fueron:
\begin{itemize}
\item \textit{grow\_policy}: \textit{depthwise}.
\item \textit{learning\_rate}: 0,1.
\item \textit{Number of estimators}: 800.
\end{itemize}
Este modelo con los hiperparámetros indicados obtuvo un \textit{F-score} de 0,69 y un \textit{G-mean} de 0,71 siendo estos los mejores valores reportados entre todos los modelos.

\section{Fase de desarrollo de la aplicación}
\label{Fase de desarrollo de la aplicación}
\subsection{Calidad el código}
Para evaluar y garantizar la calidad del código se hizo uso de la herramienta \textit{SonarCloud}. \textit{SonarCloud} es una plataforma de análisis de código y gestión de la calidad del software en la nube. Ofrece herramientas para detectar problemas en el código, tales como errores, vulnerabilidades, y código duplicado, y promueve prácticas de desarrollo de software de alta calidad.

Algunas de sus características principales son:
\begin{itemize}

\item \textbf{Análisis de Código}: Examina el código fuente para identificar errores y vulnerabilidades, asegurando que el software sea robusto y seguro.

\item \textbf{Integración Continua}: Se integra con herramientas de integración continua (CI) como Jenkins, GitHub Actions, Azure DevOps, y más, permitiendo análisis automáticos con cada commit o push al repositorio.

\item \textbf{Soporte Multilenguaje}: Soporta una amplia variedad de lenguajes de programación, incluyendo Java, JavaScript, TypeScript, C\#, C++, Python, PHP, y muchos más.

\item \textbf{Informes y Métricas}: Proporciona informes detallados y métricas sobre la calidad del código, facilitando el seguimiento del progreso y la identificación de áreas que necesitan mejora.
\end{itemize}

El uso de esta herramienta ha sido clave ya que ciertas vulnerabilidades que no se habían tenido en cuenta por desconocimiento de su solución o de las propias vulnerabilidades, fueron solucionadas de forma ágil ya que la propia herramienta muestra las soluciones y la ubicación de los errores.

\imagen{./sonarCloudReport.png}{Resumen del reporte de \textit{SonarCloud} sobre PDAIvolution.}{}

Como se observa en el reporte de \textit{SonarCloud} (ver figura~\ref{fig:./sonarCloudReport.png} existen varias categorías principales:
\begin{itemize}
\item \textit{\textbf{Security}}: muestra el número de problemas de seguridad que tiene la aplicación. Suelen tratarse de problemas como dejar visibles calves secretas, no tener en cuenta vulnerabilidades~\ldots En el caso del proyecto, no existe ningún problema de seguridad.
\item \textit{\textbf{Reliability}}: muestra el número de problemas de fiabilidad que tiene la aplicación. Se tratan de problemas como el uso de una variable antes de su declaración. En el caso del proyecto, no existe ningún problema de fiabilidad.
\item \textit{\textbf{Maintainability}}: muestra el número de problemas de mantenimiento que tiene la aplicación. En el caso del proyecto, no existe ningún problema de mantenimiento.
\item \textit{\textbf{Duplications}}: muestra el porcentaje de lineas duplicadas que tiene la aplicación. En el caso del proyecto, el porcentaje de lineas duplicadas es 0\%.
\item \textit{\textbf{Accepted Issues}}: muestra el número de problemas marcados como aceptados, es decir, problemas que se tiene conocimiento de ellos pero no se van a solucionar. En el caso del proyecto tiene marcados tres problemas como aceptados, dos de los cuales vienen dados por el proyecto PADDEL y el restante se debe a que la solución para el problema resulta más compleja que el problema en sí.
\item \textit{\textbf{Security Hotspots}}: son secciones del código que pueden no ser vulnerabilidades de seguridad por sí mismas, pero que merecen una revisión más detallada por parte de un desarrollador experimentado.
\end{itemize}

Por lo visto en el resumen, el proyecto, a excepción de los problemas previamente comentados, cumple los criterios de calidad de código.


\subsection{Problemas durante el desarrollo}

\subsubsection{Planificación temporal}

Durante el desarrollo de la aplicación, hubo dos problemas principales a la hora de realizar la parte de planificación temporal. El primero de ellos fue no conseguir conectar la aplicación de Zube con el repositorio de Github haciendo que este crease los \textit{issues} de forma automática por lo que se opto por indicar las tareas que se realizaban en cada uno de los \textit{commits}. El segundo problema se dio con el proyecto ya avanzado. Este problema fue la perdida total del proyecto en Zube. Debido a un fallo, que todavía se desconoce, el proyecto se eliminó por completo de la plataforma, eliminando a su vez todas las tareas de cada \textit{Sprint} así como la duración de los mismos. Este problema no se consiguió solucionar por lo que no se pudo recuperar nada de la información del proyecto.

\subsubsection{Compatibilidad}

Tanto durante la implementación del proyecto PADDEL como durante la fase de experimentación y despliegue, han habido problemas de compatibilidad.

Este problema de compatibilidad viene dado por la versión de CUDA~\footnote{Fuente:~\url{https://docs.nvidia.com/cuda/cuda-toolkit-release-notes/index.html}}y numba~\footnote{Fuente:~\url{https://numba.pydata.org/}}. CUDA (Compute Unified Device Architecture) es una plataforma de computación paralela y un modelo de programación creada por NVIDIA y numba es un compilador justo a tiempo (JIT) para Python que traduce funciones numéricas escritas en Python y NumPy a código máquina altamente optimizado utilizando el compilador LLVM. El problema se da cuando la versión de CUDA y numba no son compatibles entre sí. Cuando se da este problema, se ha de modificar las versiones de uno de los dos para que sean compatibles; sin embargo, al hacer la experimentación y despliegue en un servidor donde no se puede modificar la versión de CUDA, se termino resolviendo mediante la desactualización de la versión de numba provocando a su vez una desactualización de la versión de Python de la 3.12 a la 3.10. Es por este motivo por el que el proyecto no se encuentra en la versión más actualizada.

\subsection{Despliegue}

Por último, se destaca el despliegue en un servidor Linux real.
Primero se planteo el despliegue en un servicio gratuito pero, debido al gran tamaño de la aplicación y de los datos con los que se iba a trabajar, se descartó esta opción tras comprobar el reducido tamaño que ofrecen los servicios gratuitos. Tras esto, se opto por el despliegue en un servidor proporcionado por la Universidad de Burgos asegurando que la aplicación estuviera operativa y accesible de manera continua y confiable, pudiendo satisfacer las necesidades de uso. 
Durante la fase de desarrollo la aplicación fue probada en un entorno local haciendo uso de la herramienta integrada de Flask que permite lanzar un servidor de prueba y modificarlo en tiempo real. No obstante, a la hora del paso a producción, la forma en la que se despliega es mediante la herramienta de Gunicorn~\footnote{Fuente:~\url{https://gunicorn.org/}}. Gunicorn \textit{Green Unicorn} es un servidor HTTP Python WSGI para UNIX.

A su vez se creo una imagen de \textit{DockerHub}~\footnote{Fuente:~\url{https://hub.docker.com/}} (disponible en~\url{https://hub.docker.com/repository/docker/jornez/pdaivolution/general} que permite hacer un despliegue muy sencillo sin necesidad de instalar nada a excepción de \textit{Docker} o \textit{Docker Desktop}.




\capitulo{6}{Trabajos relacionados}

Durante los últimos años se han utilizado diferentes métodos para evaluar la enfermedad e Parkinson mediante la prueba del \textit{Rapid Finger-tapping}. En este capítulo se recogen algunos de los trabajos de manera resumida.

\section{A computer vision framework for finger-tapping evaluation}

En este artículo~\cite{khan2014computer} se documenta el uso de la visión por computador para la clasificación de individuos según el nivel de gravedad de la enfermedad de Parkinson.

Se emplea un método en el que se emplea el uso de la cara para la calibración de la amplitud de la prueba. Para ello el sujeto ha de elevar las manos a la altura de la cara y apuntar con las puntas de los dedos hacia la misma.

El estudio se realizó con 387 vídeos de \textit{rapid finger-tapping test} (RFT) de 13 pacientes diagnosticados con Enfermedad del Parkinson en estado avanzado y 84 vídeos de \textit{rapid finger-tapping test} de 6 personas de control sanas. En total 471 vídeos.

\subsection{Metodología}
\begin{itemize}
	\item Se realiza un reconocimiento facial del individuo y se crea a partir de el 2 cuadrados en los costados de la cara para realizar la medición de la amplitud.
	\item Se genera una serie temporal que representa la distancia desde el dedo indice al pulgar siendo esta la amplitud del movimiento.
	\item Se extraen algunas características de la serie temporal, por ejemplo, la velocidad media de apertura y cierre de dedos, el número total de toques de los dedos, la amplitud máxima...
	\item Se selección las características no redundantes mediante el algoritmo chi-cuadrado
	\item Se entrena una máquina de vectores de soporte (SVM) mediante las características obtenidas para realizar la clasificación.

\end{itemize}
\subsection{Resultados}
Se encontró una nueva característica representativa del ritmo de golpeteo llamada <<correlación cruzada entre los picos normalizados>>, la cual mostró una fuerte correlación de Guttman con las valoraciones clínicas. Al utilizar el clasificador de máquina de vectores de soporte y una validación cruzada de 10 grupos, se logró categorizar las muestras de pacientes en los niveles UPDRS-FT con una precisión del 88\%. Este mismo esquema de clasificación también permitió discriminar entre las muestras \textbf{RFT} de los controles sanos y los pacientes con \textbf{EP}, logrando una precisión del 95\%.


\section{The discerning eye of computer vision}
Este artículo~\cite{williams2020discerning} utilizan 133 vídeos de manos realizando el \textbf{RFT} procedentes de 39 pacientes de enfermedad de Parkinson y 30 pacientes de control.

A través de estos vídeos se han logrado extraer características y comprobar la relación que existe entre estas y la gravedad de la enfermedad en un paciente siendo esta valorada desde 0 (normal) hasta 4 (enfermedad muy grave).


\subsection{Metodología}
\begin{itemize}


	\item Se utiliza la librería de DeepCutLab la cual se trata de una librería de visión por computador para obtener la serie temporal de la amplitud.
	\item Se normaliza la serie con la amplitud máxima siendo igual a 1 y escalando los valores conforme a esta medida.
	\item Al igual que en el estudio anterior se extraen ciertas características de la serie siendo estas la velocidad, la amplitud y el ritmo.
\end{itemize}
\subsection{Resultados}
DeepLabCut rastreó y midió con fiabilidad el \textit{finger-tap} en un vídeo estándar de smartphone. Las medidas del ordenador se relacionaron bien con las valoraciones clínicas de la bradicinesia.


\section{Supervised classification of bradykinesia}
Este artículo~\cite{williams2020supervised} utiliza 70 vídeos de evaluaciones de \textit{finger-tap} en un entorno clínico (40 manos con \textbf{Parkinson}, 30 manos de control).
Dos expertos clínicos en \textbf{Parkinson}, que desconocían los diagnósticos, evaluaron los vídeos para dar un grado de gravedad de la bradicinesia entre 0 y 4 utilizando la Escala Unificada de Calificación de la Enfermedad de Parkinson (UPDRS, por sus siglas en inglés)


\subsection{Metodología}
\begin{itemize}
	\item Extraer la frecuencia: La frecuencia de intervención se estimó como la frecuencia correspondiente al pico de amplitud máxima en el espectro de la transformada rápida de Fourier (FFT). 
	\item La densidad espectral de energía se calculó como la integral al cuadrado del espectro FFT, una medida que se espera que aumente con la amplitud del golpeteo.
\end{itemize}
\subsection{Resultados}
Una máquina de vectores de soporte con núcleos de función de base radial predijo la presencia de bradicinesia leve/moderada/grave con una precisión de prueba estimada del 0,8 \%.
Un modelo Naïve Bayes predijo la presencia de la enfermedad de Parkinson con una precisión de prueba estimada de 0,67. 

\section{Paddel}

Este proyecto~\cite{paddelRepo} utiliza 158 vídeos de evaluación etiquetados para determinar la presencia o ausencia de Parkinson. El objetivo es realizar una predicción binaria sobre si una persona padece la enfermedad de Parkinson.

\subsection{Metodología}
\begin{itemize}


\item Extracción de «landmarks»: Convertir el vídeo de la prueba de «finger-tapping» en una secuencia temporal de puntos de referencia de la mano.

\item Lectura de imágenes del vídeo: haciendo uso de la librería OpenCV en Python para leer las imágenes del vídeo sin cargar todo el archivo en memoria.

\item Extracción de «landmarks» de las imágenes: Utilizar Mediapipe Hands para obtener puntos de referencia de la mano en cada imagen del vídeo.

\item Manejo de oclusión de «landmarks»: Descartar puntos de referencia no fiables debido a la perspectiva de los vídeos de «finger-tapping».

\item Extracción precisa de «landmarks»: Elegir no paralelizar el proceso para obtener resultados más precisos en la aplicación web.

\item Extracción de series temporales: Generar una secuencia de datos que represente la pose de la mano en cada fotograma del vídeo.

\item Extracción de características: Obtener atributos útiles para entrenar el modelo de aprendizaje a partir de las series temporales.


\end{itemize}
\subsection{Resultados}
Una máquina de vectores de soporte con los hiperparametros optimizados por medio de <<gridsearch>> da una precisión media del 78,85\%.
\capitulo{7}{Conclusiones y Líneas de trabajo futuras}

En este último aparatado, se recopilan las conclusiones y líneas de trabajo futuras del proyecto.

\subsection{Conclusiones}

Durante el desarrollo del proyecto, se ha extraído una serie de aprendizajes y conclusiones que se describen a continuación. 
\begin{itemize}


\item \textbf{La importancia de saber tratar los datos}: durante la fase de experimentación, los resultados obtenidos al principio fueron pésimos, llegando a ser peores que un clasificador aleatorio. Esto se debía al desbalanceo de las clases, el cual no se había tenido en cuenta durante las primeras fases de experimentación. Tras aplicar las técnicas de remuestreo apropiadas, se observó una mejora notable en la eficacia de los modelos.

\item \textbf{La viabilidad del uso de IA en la evaluación de la bradicinesia}: los resultados obtenidos demuestran que, mediante técnicas de visión artificial y aprendizaje supervisado, se puede llegar a evaluar de manera más o menos precisa el grado de bradicinesia del paciente. Esta precisión se puede llegar mejorar, por ejemplo, con la ampliación del \textit{dataset}.

\item \textbf{La importancia del correcto entrenamiento de modelos}: tanto en la parte de selección de características como en la parametrización de los modelos, se ha visto una gran importancia reflejada en los resultados. Se ha comprobado como un cambio en el número de características o una modificación en los parámetros pueden influir mucho a la hora de entrenar un modelo.

\item \textbf{Los resultados no siempre se obtienen resultados óptimos}: el objetivo de este proyecto no era desarrollar un modelo capaz de clasificar correctamente una alto porcentaje los casos y videos, sino aprender y descubrir con los recursos de los que se dispone cual es el limite que se puede alcanzar. 

\end{itemize}
\subsection{Líneas de trabajo futuras}

Algunas de las líneas de trabajo futuras son:

\begin{itemize}
\item Añadir un sistema de predicción integrado en la web, utilizando para ello tanto la medicina como la evolución del paciente.

\item Probar a entrenar modelos más complejos o realizar aprendizaje semisupervisado para comprobar los resultados frente al proyecto actual.
\item Buscar alguna forma para acelerar el procesado de los vídeos ya que supone una espera muy larga al utilizar la aplicación web. 

\item Internacionalizar la página web, no solo a nivel de idioma sino también a nivel cultural.

\item Una línea de trabajo a más largo plazo, es la ampliación del \textit{dataset} mediante videos de personas con enfermedad de Parkinson mostrando su evolución, como vídeos de personas sin la enfermedad para tener más muestras de grupos de control.

\end{itemize}



\bibliographystyle{plain}
\bibliography{bibliografia}

\end{document}
